\clearpage\fancyhead[L]{\small KOMET.}\hypertarget{native-408-part-lower}{}\label{native-408-part-lower}
\begin{center}\large KOMET.\end{center}
\begin{center}\hrule\vspace{4pt}\fontsize{10}{12}\selectfont Göttingische gelehrte Anzeigen. Stück 60. Seite 593..594. 13. April 1818.\par\vspace{4pt}\hrule\vspace{4pt}\end{center}
\par\indent Um die Beobachtung des Kometen zu erleichtern, den Herr \textsc{Pons} Ende vergangenen Jahres entdeckt hat, hat Herr Hofrat \textsc{Gauss} seine Bahn für die restliche Zeit seiner Sichtbarkeit vorausberechnet. Wir glauben, den Beobachtern mit der Veröffentlichung dieser Ergebnisse einen Dienst zu erweisen. Die parabolischen Elemente, nach denen diese Ephemeride berechnet wurde, beruhen auf den Beobachtungen von Herrn Dr. \textsc{Olbers} vom 3. bis 28. März und lauten wie folgt:

\begin{center}\fontsize{10}{12}\selectfont
Zeit des Durchgangs durch die Sonnennähe
\par 26. Febr. 1818. $21^{\mathrm u}\;34^{\mathrm m}\;42^{\mathrm s}$ Zeit am Meridian von Göttingen
\par\smallskip\begin{tabular}{lrr}
Kleinster Abstand von der Sonne & $\cdot\;\cdot\;\cdot\;\cdot$ & $1.20053$\\
Aufsteigender Knoten & $\cdot\;\cdot\;\cdot\;\cdot\;\cdot\;\cdot$ & $70^\circ\;5'\;12''$\\
Neigung der Bahn & $\cdot\;\cdot\;\cdot\;\cdot\;\cdot\;\cdot\;\cdot$ & $90\;0\;0$\\
Länge der Sonnennähe & $\cdot\;\cdot\;\cdot\;\cdot\;\cdot\;\cdot$ & $183\;22\;58$\\
\end{tabular}\end{center}

\SourceLeaf{409}\par\indent Erst die künftige Verbesserung dieser Elemente anhand weiterer Beobachtungen wird zeigen, ob die Neigung merklich vom rechten Winkel abweicht und in welcher Richtung. Bisher bleibt unbestimmt, ob die Bewegung in Bezug auf die Ekliptik rechtläufig oder rückläufig ist. Die Zeit der Ephemeride ist $14^{\mathrm u}\,20^{\mathrm m}$ am Meridian von Göttingen. [Die Dphemeride ist im vorhergehenden Aufsatz abgedruckt.]

\par\indent Obwohl der Komet seine Sonnennähe schon weit hinter sich hat, wird sein Licht wegen der weiterhin abnehmenden Entfernung von der Erde noch zunehmen, sodass es am 12. Mai doppelt so stark sein wird wie am 13. März. Da es aber an sich so außerordentlich schwach ist, wird der Komet nur mit guten Fernrohren zu beobachten sein.

\par\indent Die Beobachtungen dieses Kometen an der hiesigen Sternwarte werden bei einer anderen Gelegenheit veröffentlicht.
